% !TEX root = ../do-an.tex
% ---------------------------------------------------------------------
%  LỜI MỞ ĐẦU
% ---------------------------------------------------------------------
\chapter*{Lời mở đầu}
\addcontentsline{toc}{chapter}{Lời mở đầu}

Trong những năm gần đây, phong trào rèn luyện thể dục thể thao, đặc biệt là bộ môn cầu lông, đang phát triển vô cùng mạnh mẽ tại Việt Nam nói chung và thành phố Đà Nẵng nói riêng. Cầu lông thu hút đông đảo người chơi ở mọi lứa tuổi nhờ tính linh hoạt, lợi ích sức khỏe và tính gắn kết cộng đồng cao. Sự gia tăng nhanh chóng về số lượng người chơi đã thúc đẩy hàng loạt cụm sân cầu lông tư nhân và công cộng ra đời.

Tuy nhiên, công tác vận hành và đặt lịch tại đa số các cụm sân hiện nay vẫn chủ yếu thực hiện theo phương thức truyền thống: người chơi gọi điện thoại trực tiếp hoặc nhắn tin qua Zalo/Facebook cho chủ sân để hỏi giờ và đặt lịch. Quy trình thủ công này bộc lộ rất nhiều bất cập đối với cả khách hàng lẫn người quản lý:
\begin{itemize}
    \item \textbf{Đối với khách hàng:} Khó nắm bắt được tình trạng sân trống theo thời gian thực, mất nhiều thời gian chờ đợi phản hồi từ chủ sân, dễ xảy ra tình trạng đến nơi thì sân đã bị người khác sử dụng do trùng lịch.
    \item \textbf{Đối với chủ sân / nhân viên bốt trực:} Rất dễ ghi nhầm hoặc sót thông tin đặt lịch vào giờ cao điểm, tốn thời gian kiểm tra và đối soát các giao dịch chuyển khoản ngân hàng thủ công, gặp khó khăn trong việc quản lý doanh thu dịch vụ phụ trợ (nước giải khát, thuê vợt, bán cầu lông) cũng như theo dõi tỷ lệ lấp đầy sân để tối ưu hóa kinh doanh.
\end{itemize}

Nhằm giải quyết triệt để những khó khăn trên, đề tài \textbf{"Ứng dụng đặt sân cầu lông trực tuyến"} được nghiên cứu và xây dựng. Đồ án hướng tới mục tiêu cung cấp một giải pháp chuyển đổi số toàn diện: giúp người chơi dễ dàng tra cứu, theo dõi ma trận sân trống và đặt lịch thi đấu chỉ với vài lượt chạm trên thiết bị di động, tự động thanh toán qua mã VietQR thông minh; đồng thời hỗ trợ chủ sân số hóa toàn bộ quy trình điều hành bốt trực POS, kiểm soát dòng tiền minh bạch và quản lý khách hàng chuyên nghiệp.

Nội dung báo cáo đồ án được cấu trúc thành 2 chương chính:
\begin{itemize}
    \item \textbf{Chương 1: Tổng quan cơ sở lý thuyết} -- Trình bày tổng quan về các Framework cài đặt (React Native, Expo, Node.js, Express.js), danh mục các Công cụ (Tools) hỗ trợ phát triển (VS Code, Android Studio, MySQL Workbench, Draw.io, Figma, Postman) và các Ngôn ngữ lập trình cốt lõi (JavaScript, SQL, JSON).
    \item \textbf{Chương 2: Thực nghiệm} -- Trình bày chi tiết các kết quả đạt được gồm: Sơ đồ tuần tự (Sequence Diagram) cho 4 chức năng cốt lõi, Sơ đồ mạng PERT và phương pháp tính toán tìm đường găng dự án, Demo giao diện thực tế theo quy trình thiết kế 3/5 bước, và thiết kế cấu trúc Cơ sở dữ liệu quan hệ (ERD).
\end{itemize}
